\section{Introduction}
\label{sec:intro}

Differential privacy~\cite{dwork2006calibrating} is widely regarded as the
canonical formal guarantee for privacy-preserving machine learning. In its
classical form, a single privacy budget $(\eps, \delta)$ is shared by every
record, and every record receives the same protection. Yet privacy
preferences are heterogeneous: a clinical patient and a casual user
contributing data to the same model often place very different values on
their privacy. \emph{Individual differential privacy} (iDP) replaces the
single $(\eps,\delta)$ guarantee with a per-record contract
$(\eps_i,\delta)$, and several mechanisms have been proposed to deliver such
per-record contracts in deep
learning~\cite{boenisch2023have, jorgensen2015conservative, feldman2021individual}.

Heterogeneous privacy budgets have a natural home in
\emph{federated learning}~\cite{mcmahan2017communication,
kairouz2021advances}, where the data contributors are already separately
identifiable entities---silos in a hospital consortium, clients of a
mobile keyboard, banks in a fraud-detection coalition---and it is exactly
these entities that hold heterogeneous preferences about how much
protection they demand. Cross-silo federated learning with per-silo (or
per-record) personalised DP is now an active research
programme~\cite{liu2024crosssilo}, and the setting we study should be
read as its strategic component: contributors declare heterogeneous
budgets to a coordinator, the coordinator reconciles them through a
single mechanism, and the resulting per-contributor risk depends on the
budget vector of the whole population. The most widely-adopted way of
doing this reconciliation in DP-SGD is the sampling-based iDP of
\citet{boenisch2023have}, which calibrates each record's Poisson
inclusion probability $p_i$ against a single shared noise multiplier
$\sigma$ so that, for fixed batch size and iteration count, every group
of records satisfies its own $(\eps_g, \delta)$-DP bound. We build the
formal development on this mechanism because it is the empirical
canonical case, but the strategic phenomena we identify apply to any
budget-coordination scheme that jointly solves for
$(\sigma, p_g)$ from a public budget vector (\cref{sec:related}).

Because every record shares the same noise multiplier and a common
expected batch size, each record's sampling rate is mechanically coupled
to every other record's chosen budget. Lowering one's own $\eps_i$
forces other contributors' sampling rates upward to keep the expected
batch size fixed---which, as recent empirical work has
shown~\cite{kaiser2026your}, inflates their realised membership-inference
(MIA) advantage even though every formal $(\eps_g,\delta)$-iDP guarantee
continues to hold. The promise of individual privacy control is
therefore broken in a structural sense: an individual's \emph{actual}
privacy risk depends not on her own choice but on the joint choice of
the entire population.

\paragraph{Our viewpoint: this is a game.}
The natural framing for any choice problem whose outcome is determined by
the choices of multiple self-interested parties is game theory.
\citet{pejo2019together} introduced a two-player game on the privacy-accuracy
plane (the \emph{Price of Privacy}), but they were studying a mechanism in
which a player's privacy choice affected only her \emph{own} accuracy. The
novel feature of sampling-based iDP is that every player's privacy choice
affects every other player's \emph{privacy}, a far stronger and qualitatively
different coupling. We call this the \emph{privacy externality}, and the
game it gives rise to the \emph{Privacy Externality Game} (PEG).
Even without any coordination, the equilibrium of rational play is
asymmetric: at every pure NE one player is strictly worse off than
another in realised MIA advantage, and the coalition attack of
\citet{kaiser2026your} is the coordinated limit of this same
asymmetry.

\paragraph{Contributions.}
This paper makes the following contributions.
\begin{enumerate}
\item \textbf{Formalisation.} We define the $n$-player Privacy Externality
  Game on a continuous action space. Each player chooses a per-group budget
  $\eps_i \in [\eps^{\min}, \eps^{\max}]$; the sampling-based iDP mechanism
  then determines the noise multiplier $\sigma$ and the per-group sampling
  rates $p_g$, and each player's utility is the canonical scalar
  $u_i = V(\sigma) - w_i A_i(\boldsymbol{\eps})$ trading off model
  quality against the realised MIA advantage on the player's own data.
  We motivate the functional form of $V$ via the excess-empirical-risk bound
  of \citet{bassily2014private} for convex DP-ERM, noting the directional
  caveat as we use it in the (potentially non-convex) deep-learning regime.

\item \textbf{Externality lemma.} We prove that the externality is monotone:
  $\partial A_i / \partial \eps_j \leq 0$ for every pair $i \neq j$ on the
  binding feasible region of the mechanism (Lemma~\ref{lem:externality}).
  Lowering any other player's budget never reduces my realised vulnerability.

\item \textbf{Asymmetric Nash equilibrium.} Even under fully symmetric
  $2$-player play---identical group sizes and identical privacy
  weights $w_1 = w_2$---the PEG admits \emph{no} symmetric pure Nash
  equilibrium. At any diagonal profile $(\eps, \eps)$ both players
  face an identical unilateral-deviation incentive, and every pure NE
  therefore has one player carrying a strictly higher realised MIA
  advantage in exchange for a strictly higher shared model utility
  (\cref{sec:C1}). \emph{Even without any
  coordination, some player is always worse off at any Nash profile.}
  The asymmetric-role pattern persists across group-size asymmetries
  (\cref{sec:C2}), across privacy-weight ratios
  (\cref{sec:C7}), and across $n \geq 2$
  (\cref{sec:C8}).

\item \textbf{Coalition Stackelberg is the coordinated limit.} As a
  direct corollary of the externality lemma, a coalition of data
  holders whose joint utility places heavy weight on a target's harm
  best-responds by picking the participation-preserving minimum
  $\eps^{\star}$; the target's $A_t$ is then monotone in coalition
  size $k$ (\cref{cor:coalition-attack}). This recovers the
  budget-manipulation attack of \citet{kaiser2026your} as the
  coordinated limit of the asymmetric-NE picture above. In our $n=5$
  canonical setting the target's $A_t$ exceeds $0.7$ at $k=4$.

\item \textbf{Empirical evaluation.} We numerically evaluate the game on a
  range of small but representative scenarios mirroring
  \cite{kaiser2026your}. We report equilibrium budgets, realised MIA
  advantages, and the coalition-attack curve.
\end{enumerate}

\paragraph{Implication.}
The prior attack literature on iDP frames excess vulnerability as
something a malicious party inflicts. We argue the opposite: the
asymmetric NE of the previous bullet is the equilibrium behaviour of
any rational players, coalition or not, and the coalition attack of
\citet{kaiser2026your} is the coordinated limit of that same
equilibrium. Mitigation cannot proceed by detecting attackers; it must
proceed by changing the contract so the underlying game ceases to admit
harmful equilibria.

The paper is organised as follows. \Cref{sec:background} reviews
differential privacy, sampling-based iDP, MIA advantage, and basic game
theory. \Cref{sec:game} formalises the Privacy Externality Game.
\Cref{sec:analysis} proves the externality lemma, establishes the
asymmetric-NE structural claim, and derives the coalition-attack
recovery as its coordinated-limit corollary.
\Cref{sec:experiments} reports the numerical evaluation. \Cref{sec:related}
reviews related work and \cref{sec:discussion} concludes.
